\subsection{The Shape of the Four Weeks}\label{sec:shape}

The First Sunday of Advent is the Sunday that falls on 30 November or nearest to it (n.~20). It
can fall from 27 November to 3 December, so Advent has from twenty-two to twenty-eight days
before Christmas, and the fourth week is always cut short by the Vigil on 24 December. The
Missal heads each Sunday and each Ember day with the name of a Roman church, its station: St Mary
Major for the First Sunday, Holy Cross in Jerusalem for the Second, St Peter's for the Third, the
Twelve Apostles for the Fourth.

The Gospels of the Sundays move from the end of the world to the threshold of Bethlehem. The
First Sunday reads the signs before the Son of Man comes in power (Luke 21:25--33). The Second
reads John the Baptist in prison sending his disciples to Christ (Matthew 11:2--10), and the
Third reads John's answer to the priests and Levites who ask him who he is (John 1:19--28). The
Ember week turns to Mary: the Annunciation on the Wednesday (Luke 1:26--38) and the Visitation on
the Friday (Luke 1:39--47). The Saturday and the Fourth Sunday read John's preaching in the
desert (Luke 3:1--6), and the Vigil reads the angel's message to Joseph (Matthew 1:18--21). John
and Mary are the season's two figures: John points to the one who is coming, and from the
Annunciation on, Mary carries him.

Through all four weeks the Breviary reads Isaiah. In the Breviary of 1888, whose course of
lessons is the one followed here, the first lesson of the First Sunday's Matins begins the book
under this heading:

\begin{quote}
\latin{Incipit liber Isaiae Prophetae. Lectio i. Cap. 1. Visio Isaiae filii Amos, quam vidit
super Judam et Jerusalem, in diebus Oziae, Joathan, Achaz, et Ezechiae, regum Juda.}\par
{\smallskip\noindent\small ---\textit{Breviarium Romanum}, pars hiemalis (Ratisbon: Pustet,
1888), p.~133.\par}
\end{quote}

\noindent The heading reads \latin{Incipit liber Isaiae Prophetae} (gloss: here begins the book
of Isaiah the Prophet), and the lesson opens with the vision Isaiah saw concerning Judah and
Jerusalem in the days of four kings of Judah.
From there the first lesson of each day goes on through the prophet, chapter by chapter, from the
indictment of Jerusalem in the first chapters to the promises of the later ones. Under the Code
of Rubrics of 1960, Sunday Matins has a single nocturn of three lessons: two from Scripture and
the opening of a Father's homily on the day's Gospel (nn.~162\,a, 220). The 1962 Office therefore
reads two lessons of Isaiah each Sunday and then Gregory the Great on the Gospel. The Ember Days
are different. Each reads at Matins three lessons from a homily on its Gospel, St Ambrose on the
Wednesday and the Friday and Gregory on the Saturday, and no Isaiah (n.~222\,a). Isaiah is read
instead at their Masses: twice on the Wednesday, once on the Friday, and four times on the
Saturday in the longer form of its Mass. Appendix~\ref{sec:daily}
gives the Isaiah lesson for each day.

The Ember Days are the Wednesday, Friday and Saturday after the Third Sunday. In 2026 they fall on
16, 18 and 19 December. Under the 1960 Code they are ferias of the second class, as are the days
from 17 to 23 December; the other weekdays of Advent are of the third class (nn.~24--25).

The last seven days before the Vigil, from 17 December, have their own Office, with the great O
antiphons at Vespers; Section~\ref{sec:last-days} describes these days and the Vigil that ends
the season.
